%! TEX root = ../main.tex
\documentclass[main]{subfiles}
\usepackage{adjustbox}

\begin{document}
\chapter{提案手法}
\section{DCASE基礎モデル}
\subsection{音響イベント検出 (SED)}
\label{DCASE基礎モデル:音響イベント検出}
DCASE 2024では、タスクとしてSEDを扱う。
% 解説用の図を作成する
\begin{figure}[!htbp]
    \centering
    \includegraphics[width=80mm]{figures/sed.pdf}
    \caption{クリップ予測及びフレーム予測の、10秒の音声入力を対象とした出力例。全体におけるイベントの有無を判定するクリップ予測に対し、フレーム予測ではイベントの区間（Onset/Offset）も算出する。注: 入力のスペクトログラムと予測結果はあくまで構造を説明するためのイメージであり、実際の対応はない。}
    % 図の説明をより詳細にする。
    \label{fig:sed}
\end{figure}
\clearpage

本タスクでは、10秒間の音声データを「クリップ」、それをさらに短時間で分割した単位を「フレーム」と定義し、イベント種別及び各イベントの開始・終了時刻を予測する。
音声データ内（クリップ及びフレーム）では、複数のイベントが重複して発生する可能性があるため、SEDでは複数のイベント種別とその発生時刻を同時に予測することが求められる。
図\ref{fig:sed}に示すように、学習においてはクリップ単位では10秒間にどのイベントが含まれているか（イベントの有無）を、フレーム単位ではイベント種別に加えて、その開始・終了時刻（時間的境界）を学習対象とする。 \par

また、予測の粒度に応じてデータのラベル定義が異なる。時間的な詳細情報を含むものを「強ラベル」、イベントの有無のみを示すものを「弱ラベル」とし、それぞれフレームとクリップに対応する。

\end{document}